\begin{abstract}
\noindent
Thesis title: Integrating Graph Neural Networks and Statistical Dependence for Pharmacotherapy Knowledge Graphs

\medskip


This thesis investigates how structured knowledge represented by a pharmacotherapy knowledge graph can
 be combined with explicit statistical evidence derived from patient-level data.
The problem is studied through two complementary tasks operating on a shared synthetic pharmacotherapy 
benchmark.
Task~A addresses directed link prediction, using patient-level observations to evaluate candidate 
relations in the graph.
Task~B addresses patient-level clinical endpoint prediction as node classification task, 
using the graph as a shared structural scaffold while individual 
patient observations instantiate the node-level input.

Both tasks combine graph neural representations with explicit statistical dependence information.
In Task~A, a heterogeneous graph encoder is fused with candidate-specific statistical context derived 
from patient observations.
In Task~B, heterogeneous message passing is augmented with endpoint-oriented statistical evidence, 
while different graph architectures, propagation depths, and residual connections are evaluated.
HCR-derived coefficients are used as representations of empirical dependence.

On the synthetic benchmark, the final Task~A model is consistently strong under 
the official five-scenario validation protocol, and ablations show that explicit 
statistical context substantially improves link prediction over the graph representation alone.
In Task~B, the statistical branch is most beneficial when deeper message passing loses predictive 
information, whereas residual connections preserve the same information within the graph pathway. 
%consequently, the two mechanisms are not always additive on the synthetic graph.

The generality of these observations is further examined through cross-domain experiments 
on real-world datasets.
On leakage-corrected Last-FM*, progressive addition of explicit candidate-specific context 
systematically improves ranking quality over a heterogeneous graph encoder alone.
On a Hetionet-derived endpoint-prediction proxy, the statistical branch is again 
particularly valuable for deeper models.
On the H2GB IEEE-CIS-G card-fraud graph, residual connections stabilise learning and the 
statistical branch further improves fraud detection relative to residual message passing 
alone on a graph structure different from the synthetic pharmacotherapy graph.

The results support a common methodological conclusion: structured graph representations and explicit empirical dependence provide complementary sources of information for both graph reconstruction and patient-level prediction.
Their relative contribution depends on the task and on how effectively the graph neural network preserves relevant information through message passing.
The proxy experiments provide methodological stress tests, but their 
findings should be interpreted cautiously before any application to clinical settings.

\medskip
\noindent\textbf{Keywords:}
knowledge graphs,
graph neural networks,
heterogeneous graphs,
link prediction,
clinical endpoint prediction,
Hierarchical Correlation Reconstruction,
pharmacotherapy.
\end{abstract}
